\subsection{Matched structural prior and frontier comparisons}
The matched follow-up adds one neutral-prior visibility run~(G) per layout, paired with its existing category-conditioned visibility run~(S). The sensor initialization, instance frontend, fitted hypothesis bank, visibility estimator, candidate generator, execution contract and common score are shared. G and S select from their respective structural-prior scores at each acquired state. The comparison also includes an independently generated nearest reachable frontier baseline. It ranks observed-free frontier representatives by grid-path distance and uses predicted unknown-cell gain to break ties. All planners use the same measured map updates, action guard and 256-action cap including return. Episodes ending early retain their actual action counts.

The original frontier implementation unconditionally included the current cell as a distance-zero candidate. All four starts were inside known free space. Predicted unknown-cell gain remained positive at the same orientation, and all four runs alternated turns without translation. A supplementary frontier variant removes the unconditional candidate and keeps the rest of the selection rule. We report both versions and use the variant restricted to actual frontiers for the substantive system comparison.

\begin{table}[tb]
\centering
\small
\caption{Paired exploration outcomes. Frontier* is the original implementation with an unconditional current-pose candidate. Frontier restricts candidates to reachable observed-free cells adjacent to unknown cells. G and S use the same visibility forecast and differ in the structural class prior. Coverage is the fraction of reachable free cells correctly mapped as free. A95 and A100 are first actions reaching 95\% and 100\%, with dashes for thresholds unreached by the recorded stop. Object recall is weighted by reference surface area. Recall and scene F1 use a 5\,cm tolerance. The 256-action cap includes return.}
\label{tab:matched-frontier-prior}
\begin{tabular}{llrrrrrrc}
\toprule
Layout & Method & Cov. (\%) & A95 & A100 & Obj. (\%) & F1 (\%) & Actions & Return \\
\midrule
Chain A & Frontier* & 32.58 & -- & -- & 58.77 & 32.37 & 256 & Yes \\
 & Frontier & 100.00 & 76 & 140 & 80.57 & 97.82 & 202 & Yes \\
 & G & 100.00 & 70 & 154 & 80.17 & 99.17 & 256 & Yes \\
 & S & 100.00 & 70 & 154 & 80.17 & 99.17 & 256 & Yes \\
\addlinespace
Cycle A & Frontier* & 45.16 & -- & -- & 31.70 & 35.36 & 256 & Yes \\
 & Frontier & 100.00 & 56 & 80 & 74.44 & 95.39 & 136 & Yes \\
 & G & 100.00 & 67 & 133 & 79.16 & 99.46 & 256 & Yes \\
 & S & 100.00 & 67 & 133 & 75.36 & 99.38 & 256 & Yes \\
\addlinespace
Chain B & Frontier* & 32.86 & -- & -- & 57.53 & 30.22 & 256 & Yes \\
 & Frontier & 100.00 & 95 & 102 & 87.20 & 98.14 & 200 & Yes \\
 & G & 92.40 & -- & -- & 86.53 & 92.29 & 256 & Yes \\
 & S & 92.40 & -- & -- & 86.53 & 92.29 & 256 & Yes \\
\addlinespace
Cycle B & Frontier* & 43.17 & -- & -- & 43.10 & 35.19 & 256 & Yes \\
 & Frontier & 99.88 & 65 & -- & 85.59 & 98.60 & 220 & Yes \\
 & G & 100.00 & 97 & 191 & 86.47 & 98.65 & 256 & Yes \\
 & S & 100.00 & 97 & 191 & 86.47 & 98.65 & 256 & Yes \\
\bottomrule
\end{tabular}
\end{table}


G and S have identical complete action trajectories and native meshes in three layouts. Their coverage thresholds also match in Cycle~A. This pair shares its first 212 executed actions and observations. The first differing target comes from the same five candidates, and the actions first differ at 213. S reduces area-weighted object recall from 79.16\% to 75.36\% and scene F1 from 99.46\% to 99.38\%. Both priors reach full coverage in three layouts and remain at 92.40\% in Chain~B. The four pairs show a change attributable to the class prior, with no improvement in either coverage threshold or terminal object recall.

The frontier variant reaches 95\% coverage sooner than S in three layouts, including Chain~B where S does not reach it. Both planners reach full coverage in three layouts. Frontier reaches 100\% in fewer actions in the two layouts where both reach it, and it finishes Cycle~B at 99.88\%. S has higher terminal scene F1 in Chain~A and both cycle layouts, with a difference of only 0.04 percentage points in Cycle~B. Frontier raises Chain~B F1 from 92.29\% to 98.14\%. The frontier episodes end after 136--220 actions, while G and S use the full budget. Terminal F1 therefore compares each policy's map at its own stop under a shared cap. All twelve new full episodes return without simulated collision.
